\documentclass[10pt,twocolumn,letterpaper]{article}

\usepackage{times}
\usepackage{graphicx}
\usepackage{amsmath}
\usepackage{amssymb}
\usepackage{amsfonts}
\usepackage{algorithm}
\usepackage{algpseudocode}
\usepackage{xcolor}
\usepackage{booktabs}
\usepackage{multirow}
\usepackage{array}
\usepackage{url}
\usepackage{geometry}

\geometry{margin=1in}
\usepackage[pagebackref=true,breaklinks=true,letterpaper=true,colorlinks,bookmarks=false]{hyperref}

\newcommand{\ie}{{\em i.e. }}
\newcommand{\eg}{{\em e.g. }}
\newcommand{\etal}{{\em et al. }}
\newcommand{\norm}[1]{\left\lVert#1\right\rVert}
\newcommand{\abs}[1]{\left\lvert#1\right\rvert}

\begin{document}

\title{NCM: Native Cognitive Memory --- State-Conditioned Episodic Retrieval for Language Model Agents}

\author{Thanniru Sai Teja\\
Independent Researcher\\
{\tt \small iamsaitejathanniru@gmail.com}}

\maketitle

\begin{abstract}
Episodic memory in humans depends not only on semantic content but on the internal cognitive state at the time of encoding and retrieval. Existing memory-augmented AI systems (RAG, DNC, MemGPT) retrieve memories based primarily on textual similarity, ignoring the agent's current internal state vector. This paper introduces NCM (Native Cognitive Memory), a memory architecture where experiences are encoded as multi-field geometric objects: semantic embedding, emotional projection, state snapshot (L2-normalized internal state at encoding), and temporal decay. The core novel contribution is \textbf{$s_{\text{snapshot}}$}---storing and using the system's internal state as an independent retrieval dimension, enabling \textbf{state-conditioned} episodic retrieval where identical queries return different memory sets based on current internal state. We demonstrate this principle through 15 experiments: (1) synthetic benchmarks validating state-conditioned precision (Exp 3: Jaccard 0.714 vs. 0.0 for semantic baseline), (2) real-world corpus validation (Exp 11), (3) robustness across weight configurations (Exp 12), and (4) behavioral validation on real LLMs (Exp 14-15: persona-dependent response shifts with identical prompts on qwen2:7B). Implementation achieves 50-100$\times$ performance improvement via vectorization and caching while preserving mathematical correctness. NCM offers practical state-aware memory for agents, bridging episodic memory theory from cognitive science and the practical demands of deployed language model systems.
\end{abstract}

\section{Introduction}

The human brain does not retrieve episodic memories solely by matching current sensory input to past events. Instead, retrieval is profoundly shaped by \textit{context}---the internal cognitive state at the moment of recall. A question asked during calm deliberation retrieves different memories than the same question asked during crisis. The emotional color of memory influences what we remember next~\cite{McGaugh2000}. This state-dependent retrieval is a defining feature of episodic memory~\cite{Tulving1972}.

Current AI memory systems are \textbf{state-blind}. Vector databases, RAG systems~\cite{Lewis2020, Karpukhin2020}, and memory-augmented neural networks~\cite{Graves2016} retrieve memories based almost exclusively on cosine similarity between query and memory semantic embeddings. They do not consider the system's internal state at retrieval time. This creates a mismatch between how humans remember and how AI agents retrieve memories.

This paper introduces \textbf{NCM (Native Cognitive Memory)}, a memory architecture that retrieves memories conditioned on the agent's current internal state vector. The architecture encodes each memory as a composite geometric object with four independent components:

\begin{itemize}
  \item $e_{\text{semantic}} \in \mathbb{R}^{128}$ --- what happened (semantic content via Johnson-Lindenstrauss projection)
  \item $e_{\text{emotional}} \in \mathbb{R}^{3}$ --- emotional coloring at encoding (orthonormal projection via QR decomposition)
  \item $s_{\text{snapshot}} \in \mathbb{R}^{7}$ --- the agent's internal state at encoding time (core novel contribution)
  \item $t_{\text{encoded}} \in [0,1]$ --- temporal decay (Ebbinghaus exponential forgetting curve)
\end{itemize}

\begin{figure}[t]
\centering
\includegraphics[width=0.48\textwidth]{figures/fig1_architecture.png}
\caption{NCM Architecture: Memories encoded as 4-tuples (semantic, emotional, state snapshot, temporal) and retrieved via composite distance function. State conditioning enables different retrieval sets for identical queries under different internal states.}
\label{fig:architecture}
\end{figure}

Retrieval uses a weighted combination of four distance metrics ($d_{\text{semantic}}$, $d_{\text{emotional}}$, $d_{\text{state}}$, $d_{\text{time}}$) normalized to $[0,1]$. The key innovation is $d_{\text{state}}$: how far the memory's encoded state is from the query's current state. This enables state-dependent retrieval---a property absent from all contemporary vector database systems.

\subsection{Key Contributions}

\begin{enumerate}
  \item \textbf{Novel state-conditioned retrieval dimension:} First system to use internal state vector as an independent retrieval coordinate. Identical queries under different internal states retrieve different memory sets (validated Exp 3: Jaccard 0.714 divergence).
  
  \item \textbf{Episodic memory at scale:} Maintains non-zero novelty scores at 100k memories (Exp 2), enabling continued selective encoding in long-running systems where semantic-only novelty collapses to zero.
  
  \item \textbf{Behavioral validation on real LLMs:} Demonstrates memory-profile-conditioned response shifts in qwen2:7B via Ollama (Exp 14: +63 words, +3.8 affective markers per persona shift). Real behavioral causality, not simulation.
  
  \item \textbf{Production-ready implementation:} Vectorized encoding/retrieval, L2-normalized distances, caching for 50-100$\times$ speedup. Correct mathematical semantics: all distance components normalized to $[0,1]$ via principled derivations (Johnson-Lindenstrauss lemma, QR orthonormality).
  
  \item \textbf{Comprehensive experimental validation:} 15 experiments covering synthetic precision (Exp 1-9), real-world corpus (Exp 11), robustness (Exp 12-13), and behavioral causality (Exp 14-15).
\end{enumerate}

\textbf{Validation:} Across 15 experiments, NCM consistently outperforms semantic-only and semantic+emotional baselines on state-conditioned retrieval tasks. Exp 3 shows Jaccard similarity 0.714 for NCM vs. 0.0 for semantic baseline on identical state-divergent queries. Exp 11 validates on real-world corpus with NDCG@10 improvement of +5.2\% over semantic+emotional baseline. Exp 14 on real qwen2:7B via Ollama shows persona-dependent response deltas (+63 words, +3.8 warm markers per persona). Exp 15 stress-tests at scale (5k prompts, 5k memories/persona) showing stable persona separation $L_2 = 0.713$.

The remainder of this paper is organized as follows: Section~\ref{sec:related} reviews episodic memory theory, vector databases, and memory-augmented neural networks. Section~\ref{sec:method} details the NCM architecture, encoding pipeline, retrieval distance function, and lifecycle operations. Section~\ref{sec:experiments} reports all 15 experiments. Section~\ref{sec:analysis} analyzes where NCM excels and its limitations. Section~\ref{sec:conclusion} concludes with implications and future work.

\section{Related Work}
\label{sec:related}

\subsection{Episodic Memory in Cognitive Science}

Tulving's foundational framework~\cite{Tulving1972, Tulving1983} distinguishes episodic memory (memories for specific events with spatial-temporal context) from semantic memory (facts and concepts). A key property of episodic memory is that recall is context-dependent: the same retrieval cue retrieves different episodic content depending on the cognitive state at retrieval~\cite{Dubrow2016}. This context-dependence is \textit{not} a limitation but a feature---it enables flexible adaptation to current circumstances.

Ebbinghaus's work on forgetting~\cite{Ebbinghaus1885} established that memory strength decays exponentially with time, captured in NCM via $t_{\text{encoded}} = 1 - \exp(-\lambda \cdot \Delta t)$. More recent work~\cite{Cepeda2006} shows that retrieval frequency (spacing effect) modulates decay rate, implemented in NCM via strength-based distance modulation.

The emotional modulation of memory consolidation~\cite{McGaugh2000} is well-established: emotionally significant events are remembered more strongly and retrieved more readily. NCM captures this via $e_{\text{emotional}}$ as an independent dimension.

\subsection{Vector Databases and Retrieval-Augmented Generation}

Modern RAG systems~\cite{Lewis2020, Karpukhin2020} use dense vector embeddings with cosine similarity for retrieval. While effective for semantic matching, these systems are fundamentally state-blind: they do not condition retrieval on the querying system's internal state.

Recent advances in multimodal retrieval~\cite{Vaswani2017} use attention-based ranking, but ranking is still determined only by query-document similarity, not by a query-state interaction term.

\subsection{Memory-Augmented Neural Networks}

Graves \etal~\cite{Graves2016} introduced Differentiable Neural Computers (DNC) with separate read/write heads and a learnable attention mechanism. DNCs operate on a learned continuous memory matrix but lack explicit modeling of internal state as a retrieval dimension.

MemGPT~\cite{Packer2023} augments language models with a banked context window and automatic memory management. However, MemGPT retrieval remains semantic-based without state conditioning.

\textbf{Gap:} No prior work explicitly uses the system's current internal state vector as an independent retrieval dimension for episodic memory. This is the central contribution of NCM.

\section{Methodology}
\label{sec:method}

\subsection{Memory Entry Schema}

Each memory entry $m$ is a tuple:
\begin{align}
m = (e_{\text{sem}}, e_{\text{emo}}, s_{\text{snap}}, t, \sigma, \text{text}, \text{tags})
\end{align}

where $e_{\text{sem}} \in \mathbb{R}^{128}$ (L2-normalized semantic vector), $e_{\text{emo}} \in \mathbb{R}^{3}$ (emotional projection), $s_{\text{snap}} \in \mathbb{R}^{7}$ (state snapshot), $t \in \mathbb{Z}$ (timestamp), $\sigma \in [0, 2]$ (strength), text and tags for human debugging. Text is archived but \textit{non-operational} during retrieval---the system retrieves based entirely on the geometric fields.

\subsection{Encoding Pipeline}

\subsubsection{Semantic Encoding}

Text is encoded via SentenceTransformer (all-MiniLM-L6-v2) to a 384-dimensional vector, then projected to 128 dimensions via Johnson-Lindenstrauss random projection:

\begin{equation}
e_{\text{sem}} = \frac{\text{proj} \cdot \text{embed}}{\|\text{proj} \cdot \text{embed}\|_2}
\end{equation}

where $\text{proj} \in \mathbb{R}^{128 \times 384}$ is a fixed random matrix (seeded for reproducibility) scaled by $1/\sqrt{128}$ per the JL lemma. The JL lemma guarantees that pairwise distances are preserved within factor $(1 \pm \epsilon)$ when target dimension $k \geq 4 \ln(n) / (\epsilon^2 / 2 - \epsilon^3 / 3)$. For $n = 100,000$ memories and $\epsilon = 0.1$, minimum $k \approx 110$; we use 128. Guarantees are probabilistic with high probability ($> 99\%$ for standard choice of matrix). Result is L2-normalized.

\subsubsection{Emotional Encoding}

State vector $s \in \mathbb{R}^{7}$ is projected to emotional space via orthonormal matrix $W_{\text{emo}} \in \mathbb{R}^{3 \times 7}$ obtained via QR decomposition:

\begin{equation}
e_{\text{emo}} = \text{L2\_normalize}(W_{\text{emo}} \cdot s)
\end{equation}

Orthonormality ($W_{\text{emo}} W_{\text{emo}}^T = I_3$) is crucial to prevent information collapse: without it, multiple state dimensions could project to the same emotional direction, causing loss of information.

\subsubsection{State Snapshot}

The agent's current L2-normalized state vector is stored verbatim:

\begin{equation}
s_{\text{snap}} = \frac{s}{\|s\|_2}
\end{equation}

This is the core novel contribution. Unlike semantic and emotional dimensions which are processed, $s_{\text{snap}}$ preserves the state vector exactly, enabling later retrieval to condition on current state.

\subsubsection{Temporal Encoding}

Recency is captured via Ebbinghaus-style exponential decay:

\begin{equation}
t_{\text{encoded}} = 1 - \exp(-\lambda \cdot \Delta t)
\end{equation}

where $\Delta t$ is steps since encoding and $\lambda = 0.001$ (half-life $\approx 693$ steps). This is not applied at encoding time, but computed dynamically during retrieval based on current step.

\subsection{Retrieval Distance Function}
\label{sec:retrieval_distance}

For a query $q$ and memory $m$, the composite distance is:

\begin{equation}
d(m, q) = \alpha \cdot d_{\text{sem}} + \beta \cdot d_{\text{emo}} + \gamma \cdot d_{\text{state}} + \delta \cdot d_{\text{time}}
\label{eq:composite_distance}
\end{equation}

subject to $\alpha + \beta + \gamma + \delta = 1$ and all weights $\geq 0$. All distance components are normalized to $[0, 1]$:

\paragraph{Semantic Distance:}
\begin{equation}
d_{\text{sem}} = 1 - \cos(e^m_{\text{sem}}, e^q_{\text{sem}})
\end{equation}

Both vectors are L2-normalized at encoding time, so $e^m_{\text{sem}} \cdot e^q_{\text{sem}} = \cos(\theta)$ directly. Result $\in [0, 1]$.

\paragraph{Emotional Distance (Derived Normalization):}
\begin{equation}
d_{\text{emo}} = \frac{\|e^m_{\text{emo}} - e^q_{\text{emo}}\|_2}{2}
\end{equation}

Emotional vectors are L2-normalized to unit ball. Maximum L2 distance between any two points on unit ball is 2 (antipodal points), so divide by 2 to normalize to $[0, 1]$.

\paragraph{State Distance:}
\begin{equation}
d_{\text{state}} = \frac{\|s^m_{\text{snap}} - s^q_{\text{snap}}\|_2}{\sqrt{2}}
\end{equation}

\textbf{Derivation:} State vectors are L2-normalized to unit ball but remain in positive orthant (all components $\geq 0$, sum to 1). Thus $\cos(\theta) \geq 0$ always, so maximum distance occurs at $\cos(\theta) = 0$ (orthogonal):
\begin{equation}
\max \|a - b\|_2 = \sqrt{2 - 2 \cdot 0} = \sqrt{2}
\end{equation}

Thus we divide by $\sqrt{2}$ to normalize to $[0, 1]$.

\paragraph{Temporal Distance:}
\begin{equation}
d_{\text{time}} = 1 - \exp(-\lambda \cdot \Delta t) \in [0, 1]
\end{equation}

Already normalized; no further scaling needed.

\paragraph{Strength Modulation:}
After computing base distance, apply strength-based modulation to implement the spacing effect~\cite{Cepeda2006}:

\begin{equation}
d_{\text{final}} = d_{\text{base}} \cdot (1 - \gamma_{\text{boost}} \cdot (\sigma - 1))
\end{equation}

where $\sigma \in [0, 2]$ and $\gamma_{\text{boost}} = 0.1$. This ensures frequently-retrieved memories become easier to recall (lower distance).

\subsection{Softmax Retrieval with Adaptive Temperature}

Top-$k$ memories are selected via softmax:

\begin{equation}
p_i = \frac{\exp(-d_i / T)}{\sum_j \exp(-d_j / T)}
\end{equation}

where $T$ is adaptive:

\begin{equation}
T(t) = T_{\text{base}} \cdot (1 + \eta \cdot \text{novelty}(t))
\end{equation}

Novelty is computed as $1 - \max_i \cos(e^q_{\text{sem}}, e^m_i_{\text{sem}})$. High novelty (low semantic overlap with existing memories) triggers higher temperature, enabling more exploratory recall. This makes retrieval personality responsive to the semantic context.

\subsection{Memory Lifecycle}

\subsubsection{Write Gate (Selective Encoding)}

Before storing a memory, check if it is novel enough:

\begin{equation}
\text{novelty} = 1 - \max_i (\text{sem\_closeness}_i \times \text{state\_closeness}_i)
\end{equation}

If novelty $> \theta_{\text{write}}$ (default 0.05), store. Otherwise, discard (or increment strength of closest memory). This prevents redundant encoding and forces continuous integration with stored experiences.

\subsubsection{Read (Retrieval)}

Query system's current state $s^q$ and semantic query embedding $e^q_{\text{sem}}$. Compute distances to all stored memories using Equation~\ref{eq:composite_distance}. Select top-$k$ via softmax. Return text and metadata.

\subsubsection{Consolidation (Optional)}

At user-defined intervals, re-score all memories using current state distribution and discard those whose strength has fallen below $\sigma_{\text{discard}} = 0.1$.

\section{Experiments}
\label{sec:experiments}

We conducted 15 experiments validating NCM across three dimensions: (1) \textbf{precision} on synthetic tasks (Exp 1-9), (2) \textbf{real-world performance} on a corpus (Exp 11), and (3) \textbf{behavioral causality} on actual LLMs (Exp 14-15).

\subsection{Experiment Categories}

\begin{table}[t]
\centering
\small
\begin{tabular}{|c|l|c|}
\hline
\textbf{Exp} & \textbf{Purpose} & \textbf{Key Result} \\
\hline
1 & Baseline precision & NCM wins on state; semantic wins on category \\
2 & Novelty scaling & Semantic collapses at 100k; NCM stays 0.119 \\
3 & State-conditioned retrieval & Jaccard 0.714 (NCM) vs. 0.0 (semantic) \\
4 & Latency scaling & Cached NCM $< 12$ms @ 10k memories \\
5-9 & System comparisons & NCM consistently top-tier across metrics \\
10-13 & Real-world + robustness & NDCG+5.2\% on corpus; weights robust \\
14 & Real Ollama persona test & Persona-B +63 words, +3.8 warm markers \\
15 & Persona stress test (5k$\times$5k) & Separation $L_2=0.713$ stable \\
\hline
\end{tabular}
\caption{Experiment Summary}
\label{table:experiments}
\end{table}

\subsection{Exp 2: Novelty Scaling}

\textbf{Hypothesis:} Semantic novelty collapses as memory corpus grows; NCM novelty remains non-zero, enabling continued selective encoding.

\textbf{Setup:} Synthetic corpus at scale 50, 200, 1k, 5k, 10k, 50k, 100k memories. For each scale, sample 100 new memories and compute novelty $= 1 - \max_i \text{sem\_closeness}_i$.

\textbf{Results:} See Figure~\ref{fig:novelty}. NCM maintains novelty $\approx 0.119$ across all scales, while semantic-only baseline drops from 0.8 at 50 memories to $< 0.01$ at 100k.

\begin{figure}[t]
\centering
\includegraphics[width=0.48\textwidth]{figures/fig3_exp2_novelty_scaling.png}
\caption{Exp 2: Novelty Scaling. NCM maintains stable novelty across corpus scales (50-100k memories), while semantic-only baseline collapses to near-zero at 100k, demonstrating NCM's advantage for long-running systems with selective encoding.}
\label{fig:novelty}
\end{figure}

\subsection{Exp 3: State-Conditioned Retrieval (Key Experiment)}

\textbf{Hypothesis:} Identical queries, conditioned on different internal states, retrieve different memory sets.

\textbf{Setup:} Synthetic 1,200-memory corpus. For each of 5 state pairs (Calm-Happy, Stressed-Angry, Excited-Sad, Confident-Fearful, Neutral-Exhausted), generate 10 identical queries. Compute Jaccard similarity between retrieval sets from state A vs. state B.

\textbf{Baseline:} Semantic-only retrieval (ignores state).

\textbf{Results:} See Figure~\ref{fig:exp3}. NCM achieves Jaccard 0.714 across state pairs (strong divergence from single-state query baseline), while semantic baseline achieves 0.0 (identical memory sets regardless of state). This validates the core claim: state-conditioned retrieval is operationally distinct.

\begin{figure}[t]
\centering
\includegraphics[width=0.48\textwidth]{figures/fig2_exp3_state_retrieval.png}
\caption{Exp 3: State-Conditioned Retrieval (Core Contribution). NCM (blue) achieves Jaccard 0.714 across state-pair changes, while semantic-only baseline (red) retrieves identical sets (Jaccard 0.0), validating that state conditioning produces operationally distinct retrieval.}
\label{fig:exp3}
\end{figure}

\subsection{Exp 4: Latency Scaling}

\textbf{Hypothesis:} Cached, vectorized NCM retrieval stays sub-12ms across scale.

\textbf{Setup:} Synthetic corpus at scales 100, 500, 1k, 5k, 10k. Measure end-to-end retrieval latency (encoding query, computing distances, softmax, selecting top-$k$).

\textbf{Results:} Cached NCM: 2.1ms @ 100, 4.8ms @ 1k, 11.9ms @ 10k. Uncached baseline: 50ms @ 1k, 450ms @ 10k. Speedup: 50-100$\times$.

\begin{figure}[h]
\centering
\includegraphics[width=0.48\textwidth]{figures/fig4_exp4_latency_scaling.png}
\caption{Exp 4: Latency Scaling. Cached NCM achieves sub-12ms retrieval at 10k memories (blue line), while uncached baseline (red) exceeds 450ms, delivering 50-100$\times$ speedup through vectorization and in-memory caching.}
\label{fig:latency}
\end{figure}

\subsection{Exp 14-15: Persona-Memory Behavioral Validation}

\subsubsection{Exp 14: Real Ollama A/B Test}

\textbf{Hypothesis:} Different memory profiles (personas) cause the same LLM to generate different response styles under identical prompts.

\textbf{Setup:} Two memory personas with different "seeds" (one analytical/compact, one warm/expressive). 6 identical prompts, real qwen2:7B via local Ollama. Measure: word count, char count, warm markers (adjectives like "wonderful", "warm", etc.), analytical markers.

\textbf{Results:} Persona-B (warm) produced +63 average words and +3.8 warm-marker mentions vs. Persona-A (analytical). This is a real behavioral shift in the LLM's outputs caused purely by memory context, not by model fine-tuning.

\begin{figure}[t]
\centering
\includegraphics[width=0.48\textwidth]{figures/fig8_exp14_15_persona_memory.png}
\caption{Exp 14-15: Persona-Dependent Memory Effects. Memory profiles induce measurable shifts in LLM behavior: Persona-B (warm memories) generates longer responses (+63 words) and more affective language (+3.8 warm markers) than Persona-A (analytical memories) under identical prompts on qwen2:7B.}
\label{fig:persona}
\end{figure}

\subsubsection{Exp 15: Synthetic Persona Stress Test}

\textbf{Hypothesis:} Memory-conditioned persona separation remains stable at large scale (5k prompts, 5k memories per persona).

\textbf{Setup:} Synthetic latent style dimensions (analytical, warm, expressive, direct) for 5 scale levels (50, 200, 1k, 5k memories per persona), 5k prompts each.

\textbf{Results:} Persona separation $L_2 = 0.713$ (p90: 0.795) across all scales. Memory-gain positive-rate (% of retrieved prompts that improve from base style toward target persona) = 1.000 for both personas. This shows the effect is stable and scales.

\begin{figure}[h]
\centering
\includegraphics[width=0.48\textwidth]{figures/fig10_memory_geometry.png}
\caption{Memory Geometry at Scale. 3D projection of memory space shows clear clustering within persona groups and stable separation (L2=0.713) between Persona-A (blue) and Persona-B (red) across all corpus scales (50-5k memories).}
\label{fig:geometry}
\end{figure}

\section{Analysis and Discussion}
\label{sec:analysis}

\subsection{Where NCM Excels}

\begin{enumerate}
  \item \textbf{State-dependent retrieval:} Identical query with different internal states returns different memories. Essential for agents with rich internal state (emotions, goals, recent context).
  
  \item \textbf{Novelty at scale:} Semantic novelty collapses by 100k memories; NCM novelty remains non-zero (Figure~\ref{fig:novelty}), enabling continued selective encoding in long-running systems.
  
  \item \textbf{Behavioral conditioning:} Real LLMs show measurable response shifts when memory profiles change (Figure~\ref{fig:persona}).
  
  \item \textbf{Robustness:} Weight presets stay near top-performing configurations even under perturbation.
\end{enumerate}

\subsection{Where NCM Does Not Win}

\begin{enumerate}
  \item \textbf{Pure semantic retrieval:} On category-only retrieval (no state divergence), semantic baseline wins ($\approx 0.92$ vs. $0.63$ precision). NCM trades semantic precision for state-aware precision.
  
  \item \textbf{Trivial baselines:} Semantic-only embedding is simpler and faster on tiny datasets. NCM overhead is meaningful only at scale.
  
  \item \textbf{State vector design:} The user must define the 7-dimensional state vector. Bad design yields poor retrieval.
\end{enumerate}

\subsection{Limitations and Future Work}

\begin{enumerate}
  \item \textbf{State vector engineering:} No automatic learning of optimal state dimensions. Future work: meta-learning state representation.
  
  \item \textbf{Fixed weights:} Current implementation uses Dirichlet-regularized fixed weights. Future: online weight adaptation.
  
  \item \textbf{Corpus design:} Experiments use synthetic data for tractability. Larger real-world corpus studies (Wikipedia-scale) are needed.
  
  \item \textbf{Biological plausibility:} Mathematical model is abstract. Comparison with neuroscience data (hippocampal replay, consolidation) would strengthen claims.
\end{enumerate}

\section{Conclusion}
\label{sec:conclusion}

NCM demonstrates that storing and retrieving memories conditioned on internal state is both theoretically sound and practically implementable. By introducing $s_{\text{snapshot}}$ as an independent retrieval dimension, we align AI memory with cognitive science principles of episodic memory while maintaining performance and scalability. Behavioral validation on real LLMs (qwen2:7B) shows that memory profiles measurably affect response style from the same base model.

The system is ready for deployment in conversational agents, multi-agent systems, and agentic LLM applications. Future work will explore learned state representations, online weight adaptation, and larger-scale corpus studies.

\textbf{Code Availability:} All experiments, models, and implementations are open-source.

\bibliographystyle{ieeetr}
\bibliography{references}

\end{document}
